\scanpage{40}
\mspage{14}

\uline{Cas des surfaces conformes à bord} (pas néc. cptes)

On utilise l'équivalence de catégories
\[
\text{surfaces conformes à bord}\quad\overset{\sim}{\simeq}\quad
\begin{gathered}\text{surfaces holomorphes sans bord,}\\
\text{munies d'une antiinvolution.}\end{gathered}
\]
\diagram{\begin{tikzcd}[ampersand replacement=\&,row sep=small]
X \arrow[r,maps to] \& (X'=X^{\mathrm{or}},\sigma)\\
X=X'/\sigma \& (X',\sigma)\arrow[l,maps to]
\end{tikzcd}}

On dit que $X$ est hyperbolique sss $X^{\mathrm{or}}$ est hyperbolique \inserted{(i.e ses comp. connexes le sont)}

Quand $X$ est sans bord, on a $\widetilde X\simeq\widetilde{X^{\mathrm{or}}}$ (car $X^{\mathrm{or}}$ est un rev. étale de $X$) donc $X$ est hyperbolique au sens des réflexions précédentes, sss il l'est au nouveau sens.

Si $X$ est compact, \inserted{donc $X^{\mathrm{or}}$ aussi,} comme on a
\[
\chi(X^{\mathrm{or}})=2\chi(X)
\]
et que $\chi(X^{\mathrm{or}})<0\Longleftrightarrow X^{\mathrm{or}}\text{ hyperbolique}$, on trouve
\[
X\text{ hyperbolique}\Longleftrightarrow\chi(X)<0
\]

P.ex. si $X$ est holomorphe 0-connexe, de type $g,\unsure{\nu|v}$, cela signifie que $2g-2+\nu\geq1$ i.e $2g+\nu\geq3$

Mais on se gardera de confondre cette condition avec celle que $\mathrm{Int}\,X$ soit hyperbolique, ce qui est tjrs vrai sauf si $\nu=0$, $g\in\{0,1\}$..
% archive sheet number: 42
